% ===================================================================
%  HOTO Na 14 — Titi. --- Masu sana’a.
%
%  Traduction, faite en 2026, en haoussa d'aujourd'hui, sur l'ido de
%  texto/io. Regles de la colonne : voir 10-hoto-01.tex. Une seule
%  numerotation. C'est le catalogue des metiers, et il compte
%  quatre-vingt-dix-sept renvois.
%
%  LE HAOUSSA A DEUX FACONS DE NOMMER UN METIER, ET CE TABLEAU MONTRE
%  QUE LA DIFFERENCE EST UNE DATE. La premiere est le prefixe « ma- »
%  du tableau 1, soude au radical : maƙeri le forgeron, magini le
%  macon, manomi le fermier, makiyayi le berger, maroƙi le mendiant,
%  mahauci le boucher, makaɗi le tambourinaire, maɗinki le tailleur
%  (62), marini le teinturier (53), magatakarda le notaire (84),
%  maƙerin zinare l'orfevre (37). La seconde est le mot LIBRE « mai »,
%  celui qui a ou qui fait, pose devant n'importe quoi : mai ɗaukar
%  hoto le photographe (20), mai sayar da kayan kiɗa le marchand de
%  musique (10), mai gyaran laima le raccommodeur de parapluies (51),
%  mai kunna fitilu l'allumeur de reverberes (87), mai canjin kuɗi le
%  changeur (29).
%
%  UN METIER A NOM EN « ma- » EST PLUS VIEUX QUE LE SIECLE ; UN
%  METIER A NOM EN « mai » A ETE FABRIQUE HIER. Le prefixe est ferme :
%  on ne peut plus en faire de nouveaux, et la liste des ma- est
%  exactement celle des metiers que la ville haoussa avait avant les
%  Britanniques. Le mot libre est ouvert : il prend n'importe quel
%  verbe et n'importe quel objet, et il a nomme tout ce qui est
%  arrive depuis. Une langue qui date ses metiers par leur
%  morphologie.
%
%  ET LE COIFFEUR (23) LE PROUVE A LUI SEUL. Au tableau 7, le barbier
%  du village etait « wanzami », un mot plein, un metier hereditaire,
%  celui qui rase, coupe, circoncit et saigne. Celui-ci a une
%  boutique au premier etage, un garcon, une clientele de ville : il
%  est « mai gyaran gashi », celui qui arrange les cheveux. Les memes
%  ciseaux, le meme geste, et deux mots que trois siecles separent.
%
%  DIXIEME REFUS DE LA COLONNE, ET IL EST L'EXACT PENDANT D'UNE
%  ACCEPTATION. Le chanteur ambulant (60) tourne la manivelle d'un
%  orgue de Barbarie et ramasse des sous. Le haoussa a « maroƙi » —
%  celui qui va de maison en maison chanter la louange et recevoir un
%  don —, et la colonne l'a ECRIT au tableau 3, pour le mendiant
%  appuye sur sa bequille : la, l'homme demandait, et maroƙi veut
%  dire demander. Ici l'homme joue d'un instrument a manivelle pour
%  un public qui passe, ce qui n'est pas le metier du maroƙi mais
%  celui d'un musicien de rue. On ecrit « mawaƙin titi ». LE MEME MOT
%  ACCEPTE A UN TABLEAU ET REFUSE A UN AUTRE, POUR LA MEME RAISON LES
%  DEUX FOIS : ce que l'homme fait.
%
%  ET « keke » ATTEINT DIX DANS CE SEUL TABLEAU. Les fiacres (52)
%  sont des « kekunan doki », et les machines a coudre des couturieres
%  (63) sont des « kekunan ɗinki » — ce dernier etant precisement le
%  mot cite dans l'en-tete du tableau 2 pour montrer que keke n'etait
%  pas un emprunt fait pour la bicyclette. Le voila dans le texte, au
%  quatorzieme tableau, a sa place et sans qu'on l'ait cherche.
%
%  ENFIN L'IMPRIMERIE (75) SE DIT AVEC LE VERBE DU TAMBOUR. « buga »
%  est frapper : bugun ganga battre le tambour, bugun kari marquer la
%  mesure au tableau 1, bugun fiyano jouer du piano — et bugu
%  imprimer. Neuvieme paire d'homonymes de la colonne, et la plus
%  juste de toutes : une presse frappe une feuille, et le haoussa
%  n'a pas cru bon de le dire autrement.
% ===================================================================

%%K t14-apar-1 apar
\VUsaut{4.0mm}
\VUtitre{15.0pt}{60}{HOTO Na 14}
\VUsaut{3.0mm}

%%K t14-tit-1 sub
\VUpk{11.4pt}{40}{Titi. --- Masu sana’a.}
\VUsaut{3.0mm}

%%K t14-01-1 p
1. --- Abokina Ioannes, wanda yake zaune a ƙauye, ya zo ya yi kwana
uku a iyalina. \VUgras{Har wancan lokaci} bai sami damar ganin babban
gari ba ; shi ya sa muke amfani da dukan kwanakinmu muna yawo, ni kuwa
ina bayyana masa abin da bai sani ba.

%%K t14-02-1 p
2. --- Da farko mun je mun ziyarci \VUgras{wurin lura da taurari}
\textsuperscript{(1)}, wanda ya ba shi sha’awa ƙwarai. Muna dawowa ta
\VUgras{titi} \textsuperscript{(3)} inda \VUgras{wankan Turkiyya}
\textsuperscript{(2)} yake, sai muka haɗu da \VUgras{jana’iza}
\textsuperscript{(4)}. A gaban \VUgras{tawagar jana’izar} akwai
\VUgras{malaman coci} suna tafiya, a gaban \VUgras{motar gawa,} inda aka
ajiye \VUgras{akwatin gawa.} Abokai da yawa sun raka iyalin \VUgras{masu
makoki.}

%%K t14-02-2 p
Bayan tawagar ta wuce, muka ƙetare titin gaban babban \VUgras{gidan
giya} \textsuperscript{(5)} inda \VUgras{masu sha} da yawa suke shan
\VUgras{giya} a kan dandali, \VUgras{rufin gilashi}
\textsuperscript{(6)} yana kare su.

%%K t14-03-1 p
3. --- Ioannes ya yi mamaki ƙwarai da \VUgras{garkuwa} da aka ajiye
tsakanin shagon \VUgras{mai sayar da giya} \textsuperscript{(7)} da na
\VUgras{mai sayar da kayan kiɗa} \textsuperscript{(10)}. Ya tambaye ni
ma’anarta ; na amsa masa cewa tana nuna \VUgras{ofishin jakadanci}
\textsuperscript{(8)} na Ingila, na kuma sa ya lura da \VUgras{jakada}
\textsuperscript{(9)} wanda yake kan baranda.

%%K t14-tit-2 sub
\VUpk{10.2pt}{40}{Duba shagunna.}
\VUsaut{3.0mm}

%%K t14-04-1 p
4. --- Ba shakka mun tsaya gaban wurin baje \VUgras{kayan kiɗa,} da
\VUgras{harmoni,} da \VUgras{organ} da \VUgras{fiyano} ; muka kalli
murfin \VUgras{littattafan waƙa} masu hoto da na kiɗan fiyano, muka
kuma yaba wa \VUgras{kayan kiɗa mai zare} na katakon zinare, wanda ake
amfani da shi a matsayin alama.

%%K t14-05-1 p
5. --- Gizagizan ƙura da \VUgras{masu share titi} \textsuperscript{(11)} da
\VUgras{motar share titi} \textsuperscript{(12)} suka tayar ya tilasta mana
mu wuce da sauri gaban kyawawan \VUgras{kayan gida} na \VUgras{mai sayar
da kayan gida} \textsuperscript{(13)} da gaban wurin baje \VUgras{murhuna}
da \VUgras{fitilu} na \VUgras{mai aikin kwano} \textsuperscript{(14)}.

%%K t14-06-1 p
6. --- Ban da haka, a wannan wuri, an tilasta mana mu bar gefen
hanyar, domin ɗamummuka sun cika shi, waɗanda aka sauke daga
\VUgras{motar jigila} \textsuperscript{(16)} domin a ajiye su a cikin
\VUgras{shago} \textsuperscript{(15)} da aka yi hayarsa yanzu.

%%K t14-06-2 p
Wannan ya hana mu ganin kyawawan motocin \VUgras{mai yin motoci}
\textsuperscript{(17)}, amma bai hana mu tsayawa mu kalli \VUgras{mai yin
akwati} \textsuperscript{(19)} yana yin akwati domin maƙwabcinsa,
\VUgras{mai sayar da kekuna da motoci} \textsuperscript{(18)}. Sa’an nan,
tun da ya riga ya yi ɗan makara, sai muka koma gida.

%%K t14-tit-3 sub
\VUpk{10.2pt}{40}{Yawo cikin gari.}
\VUsaut{3.0mm}

%%K t14-07-1 p
7. --- Washegari, mahaifiyata ta ba Ioannes shawara a ɗauki hotonsa,
sa’an nan ta danƙa mini in raka shi wurin \VUgras{mai ɗaukar hoto}
\textsuperscript{(20)}. Bayan abincin rana, muka shiga \VUgras{ɗakin
aikin} mai fasahar, inda aka tilasta mana jira mai tsawo sosai. Domin
mu shagaltar da kanmu, sai muka kalli \VUgras{hotuna}
\textsuperscript{(22)} da suke rataye a bango.

%%K t14-07-2 p
Lokacin da lokacinmu ya yi, aikin bai daɗe ba, kuma nan take da
abokina ya gama tsayawa gaban \VUgras{kyamara} \textsuperscript{(21)}, sai
muka sauka.

%%K t14-08-1 p
8. --- A kan bene na farko, muka gani ta ƙofar \VUgras{mai gyaran
gashi} \textsuperscript{(23)}, wadda take a buɗe, ɗan aikinsa yana yanke
gashin wani abokin ciniki.

%%K t14-08-2 p
Da muka iso titin, muka wuce gaban \VUgras{shagon taba}
\textsuperscript{(24)}. Ioannes ya ji daɗin \VUgras{fitilar hannu} ja
\textsuperscript{(25)} da \VUgras{babban bututun taba} da ake amfani da
shi a matsayin \VUgras{alama} \textsuperscript{(26)} ƙwarai, har ya yi
karo da wasu tsofaffin malamai biyu waɗanda suke hira suna \VUgras{shan
taba ta hanci} \textsuperscript{(27)}.

%%K t14-tit-4 sub
\VUpk{10.2pt}{40}{Shagon kek.}
\VUsaut{3.0mm}

%%K t14-09-1 p
9. --- \VUgras{Takardun kuɗi} da \VUgras{tsabar kuɗi} na kowace ƙasa
da aka baje a \VUgras{tagar shago} ta \VUgras{mai canjin kuɗi}
\textsuperscript{(29)} sun ja hankalinmu na ɗan lokaci, amma tun da
lokacin ƙaramin abinci ne, sai muka shiga wurin \VUgras{mai yin kek}
\textsuperscript{(31)} ba tare da ƙara tsayawa ba.

%%K t14-10-1 p
10. --- Muna ganin dukan \VUgras{kayan zaki} da shagon yake ɗauke da
su, wato \VUgras{kek,} da \VUgras{kek na zuma,} da \VUgras{biskit} da
\VUgras{alewa} iri iri, sai muka kasa zaɓi cikin sauƙi. A ƙarshe
Ioannes ya zaɓi \VUgras{kek na kirim,} ni kuwa na ɗauki \VUgras{ƙaramin
kek na ceri} kaɗai, domin kayan zaki suna \VUgras{ba ni ciwon haƙori,}
ni kuma ba na son zuwa wurin \VUgras{likitan haƙori}
\textsuperscript{(30)} sau da yawa ko kaɗan.

%%K t14-tit-5 sub
\VUpk{10.2pt}{40}{Buga rubutu da injin.}
\VUsaut{3.0mm}

%%K t14-11-1 p
11. --- Domin in nuna wa abokina \VUgras{injin buga rubutu} wanda ya ji
labarinsa, amma wanda bai sani ba, sai muka shiga \VUgras{ofishin}
\textsuperscript{(32)} \VUgras{ɗan baffana} \textsuperscript{(34)}. Muka same
shi yana rubuta wasiƙunsa ga \VUgras{sakatarensa} \textsuperscript{(35)}
wanda yake \VUgras{mai gaggauta rubutu.} Da ƙyar wasiƙa ta ƙare,
\VUgras{mai buga rubutu} \textsuperscript{(33)} ya sake kwafinta da injinsa.
Tun da ba ma son katse aikin ɗan uwana, sai muka yi ɗan lokaci kaɗai a
wurinsa.

%%K t14-12-1 p
12. --- A ƙarƙashin ofishin ɗan baffana, akwai babban \VUgras{mai sayar
da agogo da kayan ado} \textsuperscript{(37)} wanda yake kuma
\VUgras{maƙerin zinare.} Shagonsa mai kyau yana ɗauke da
\VUgras{agogo,} da \VUgras{agogo masu lilo,} da \VUgras{agogon aljihu,}
da \VUgras{sarƙoƙi} da \VUgras{kayan ado} masu tsada, alal misali
\VUgras{sarƙar lu’ulu’u,} da \VUgras{munduwar} zinare, da \VUgras{ƴan
kunne} da \VUgras{zobba} da aka yi musu ado da \VUgras{duwatsu masu
daraja} da \VUgras{lu’ulu’u masu ƙyalli.} An keɓe ɗaya daga cikin
tagogin domin \VUgras{kayan zinare} kuma yana ɗauke da babban tarin
\VUgras{kayan ci} na azurfa.

%%K t14-13-1 p
13. --- Muna juyawa a kusurwar titin, na lura cewa an canja
\VUgras{allon suna} \textsuperscript{(36)} da aka ajiye kusa da
\VUgras{lambar} gidan. Zan faɗi abin da na lura da babbar murya, sai
aka ji sautin kiɗan \VUgras{sojoji} masu farin ciki. Muka fara gudu don
haɗuwa da rundunar, ba tare da mun yi tunani cewa \VUgras{ita} za ta
wuce gaban shagunan \VUgras{mai sayar da huluna} \textsuperscript{(39)} da
na \VUgras{mai sayar da safar hannu} \textsuperscript{(40)} ba, kuma cewa
da mun kasance a wuri mai kyau haka domin mu ga faretinta idan muka
tsaya gaban tagar \VUgras{mai sayar da tabarau} \textsuperscript{(38)},
wanda yake cike da \VUgras{tabarau mai maƙale,} da \VUgras{tabarau mai
hannu} da \VUgras{ƙananan madubin gani.}

%%K t14-tit-6 sub
\VUpk{10.2pt}{40}{Faretin tafiya.}
\VUsaut{3.0mm}

%%K t14-14-1 p
14. --- A gaban \VUgras{rundunar sojoji} \textsuperscript{(41)} akwai
\VUgras{shugaban makaɗa} \textsuperscript{(42)} yana tafiya, \VUgras{masu
ganga} \textsuperscript{(43)} suka bi shi, da \VUgras{masu busa}
\textsuperscript{(44)} da dukan \VUgras{ƙungiyar makaɗa.}
\VUgras{Janar} \textsuperscript{(45)}, wanda yake kan filin tare da
mataimakinsa, ya tsaya kusa da \VUgras{bakin ruwa}
\textsuperscript{(49)} na \VUgras{marmaro} \textsuperscript{(48)}, domin ya
halarci \VUgras{faretin.}

%%K t14-14-2 p
\VUgras{Manyan jami’ai} \textsuperscript{(46)}, da \VUgras{kanar} da
\VUgras{mataimakin kanar} sun gaishe shi da takobi. \VUgras{Manjo} yana
gaban \VUgras{bataliyoninsa} \textsuperscript{(47)} kuma kowane
\VUgras{kyaftin} yana jagorantar \VUgras{kamfaninsa,} alhali
\VUgras{laftanoni,} da \VUgras{ƙananan laftanoni} da \VUgras{sajanoni}
suka riƙe wurarensu na saba kusa da mutanensu.

%%K t14-15-1 p
15. --- Masu wucewa da yawa sun taru a kan \VUgras{mahaɗar hanyoyi}
\textsuperscript{(50)} ; ƴan kasuwa, da \VUgras{mai gyaran laima}
\textsuperscript{(51)}, da \VUgras{marini} \textsuperscript{(53)}, da
\VUgras{maƙerin makamai} \textsuperscript{(54)} sun fito domin su ga
\VUgras{sojoji.} Dukan abin hawa, wato \VUgras{kekunan doki}
\textsuperscript{(52)}, da \VUgras{motoci masu mastu}
\textsuperscript{(57)} da kuma \VUgras{motar shayar da titi}
\textsuperscript{(56)}, sun ajiye kansu a gefen hanyar.

%%K t14-tit-7 sub
\VUpk{10.2pt}{40}{Fage a bakin titi.}
\VUsaut{3.0mm}

%%K t14-15-2 p
\VUgras{Ɗan kasuwa mai yawo} \textsuperscript{(55)} da yake ɗauke da
\VUgras{akwatunan samfurinsa} a hannu, ya ƙetare titin da sauri,
mutanen da suke zaune a kan \VUgras{benen sama} \textsuperscript{(59)} na
\VUgras{babbar motar haya} \textsuperscript{(58)} kuwa suka juya domin su ji
daɗin abin kallon. \VUgras{Mawaƙin titi} \textsuperscript{(60)} mai tausayi
da \VUgras{injin waƙarsa mai juyawa} \textsuperscript{(61)} ya tilasta
katse \VUgras{waƙarsa,} domin hayaniyar ganguna ta rufe muryarsa.

%%K t14-16-1 p
16. --- Lokacin da rundunar ta iso gaban \VUgras{shagon yadi}
\textsuperscript{(64)}, \VUgras{mata masu ɗinki} \textsuperscript{(63)} da
\VUgras{maɗinka} \textsuperscript{(62)} sun bar \VUgras{kekunan ɗinkinsu} da
aikinsu domin su kalla ta taga. \VUgras{Ɗan kasuwa}
\textsuperscript{(65)}, wanda ya ajiye sabbin \VUgras{kaya}
\textsuperscript{(66)} a \VUgras{wurin baje kayansa} yanzu, ya tilasta sa
\VUgras{yadudduka} \textsuperscript{(67)} da zannuwan da aka kafa a gefen
hanya, kusa da \VUgras{bakin magudanar ruwa} \textsuperscript{(68)}, a
ciki, don tsoron kada taron ya kife su ; har tsoho \VUgras{mai sayarwa
mai yawo} \textsuperscript{(69)}, wanda wata mai gadin ƙofa take ciniki da
shi a kan safa, an tilasta masa ya shiga wata hanya domin ya gama
sayarwarsa.

%%K t14-16-2 p
Mun raka rundunar har bariki, \VUgras{kuma,} muna gamsuwa ƙwarai da
yininmu, muka komo a lokacin abincin dare.

%%K t14-tit-8 sub
\VUpk{10.2pt}{40}{Sana’o’i iri iri.}
\VUsaut{3.0mm}

%%K t14-17-1 p
17. --- Washegari, mahaifina ya ba mu shawara mu ziyarci gidan bugu na
jaridar wurin. Mun karɓa da farin ciki.

%%K t14-17-2 p
\VUgras{Mai bugu} shi ne kuma \VUgras{mai sayar da littattafai}
\textsuperscript{(70)}, da \VUgras{mai fitar da littattafai,} da
\VUgras{mai sayar da takarda} da \VUgras{mai ɗaure littattafai.} Ya kafa
a bene na farko \VUgras{ofishin edita} \textsuperscript{(71)} na jaridar,
da kuma \VUgras{ɗakin zana,} inda \VUgras{masu zana a kan dutse}
\textsuperscript{(72)} da \VUgras{masu zana a kan jan ƙarfe}
\textsuperscript{(73)} suke aiki, a ƙarshe kuma \VUgras{ɗakin tsara
haruffa,} inda \VUgras{ma’aikatan bugu} \textsuperscript{(74)} suke.
\VUgras{Wurin bugu} \textsuperscript{(75)} da kansa, da \VUgras{injinan
bugu} da injina iri iri suna a benen ƙasa.

%%K t14-tit-9 sub
\VUpk{10.2pt}{40}{Ioannes yana tambaya da yawa ƙwarai.}
\VUsaut{3.0mm}

%%K t14-18-1 p
18. --- Muna fitowa daga gidan bugun, Ioannes ya tsaya cikin mamaki
ƙwarai gaban ƙaramin akwati na gilashi da aka ajiye a kan ginshiƙi, a
gaban \VUgras{shagon ƙananan kaya} \textsuperscript{(77)}, ya kuma tambayi
me ake amfani da shi. Mahaifina ya bayyana masa cewa \VUgras{na’urar
sanar da gobara} \textsuperscript{(76)} ce. Ban da haka kome, a cikin
garin, abin mamaki ne ga abokina, tun daga \VUgras{marmaron dutse}
\textsuperscript{(78)} mai sauƙi da \VUgras{keke mai ƙafa uku na kaya}
\textsuperscript{(80)} mai sauƙi da \VUgras{ɗan aike}
\textsuperscript{(79)} sanye da kayan aiki yake tuƙa shi, har zuwa
\VUgras{motar waya} \textsuperscript{(81)}.

%%K t14-19-1 p
19. --- Alhali mahaifina ya bar mu na ɗan lokaci domin ya yi magana da
\VUgras{mai karɓar haraji} \textsuperscript{(83)}, wanda ya tsaya domin ya
yi hira da \VUgras{lauya} \textsuperscript{(82)} da \VUgras{magatakarda}
\textsuperscript{(84)}, sai muka shagaltar da kanmu muna bin zirga-zirga a
kan titin da idanunmu. Mutane iri iri sun wuce a gabanmu :
\VUgras{ɗan aiken mai yin kek} \textsuperscript{(85)}, yana ɗauke da
\VUgras{kek mai daɗi} mai kyau cikin kwando, ya biyo bayan \VUgras{mai
sayar da tufafi} \textsuperscript{(86)} ; \VUgras{mai kunna fitilu}
\textsuperscript{(87)} yana hira da \VUgras{ma’aikacin gas}
\textsuperscript{(88)}, \VUgras{ƴar’uwar addini} \textsuperscript{(89)} wadda
take riƙe da hannun wata maraya tana komawa gidan ƴan’uwan cocinta.

%%K t14-tit-10 sub
\VUpk{10.2pt}{40}{A lokacin dawowarmu gida.}
\VUsaut{3.0mm}

%%K t14-20-1 p
20. --- A daidai lokacin da mahaifina ya haɗu da mu, Ioannes ya yi
dariya ƙwarai, yana ganin fuska mai ƙazanta da tufafi masu ban mamaki
na \VUgras{masu tsaftace bututun hayaƙi} \textsuperscript{(91)} biyu cike
da kirin.

%%K t14-21-1 p
21. --- Na so in saya wurin \VUgras{mai sayar da furanni}
\textsuperscript{(92)} tsiro domin mahaifiyata, amma mahaifina ya sa na
lura cewa zai yi nauyi ƙwarai a ɗauka kuma cewa ya fi kyau a sayi
\VUgras{lemu.} Muka kusanci \VUgras{mai sayarwa} \textsuperscript{(94)}
kuma lokacin da \VUgras{malamar gida} \textsuperscript{(93)}, wadda take
zaɓi domin ɗalibarta, ta gama sayenta, sai muka sa aka yi mana hidima.

%%K t14-22-1 p
22. --- Muna dawowa gida, muka haɗu da waɗansu sanannu : tsohuwar
abokiyar iyalina, wadda ta jingina hannunta a kan \VUgras{abokiyar
zamanta} \textsuperscript{(95)}, da \VUgras{injiniyan} \textsuperscript{(96)}
gadoji da hanyoyi, da ɗaya daga cikin ƴan goggona tare da \VUgras{mai
kula da yaranta} \textsuperscript{(98)}.

%%K t14-22-2 p
Sa’an nan muka haɗu da \VUgras{mai yin huluna} \textsuperscript{(97)} na
mahaifiyata da \VUgras{malaman addini} \textsuperscript{(99)} biyu baƙi ga
garin, waɗanda suka zo wurinmu domin su tambayi mahaifina hanyar zuwa
tasha.

\VUsaut{6.0mm}
\VUfilet{22mm}
